\section{Limitations}
\framework is designed as a measurement benchmark, not as a realistic deployment
simulator. The symbolic environments are deliberately small and cannot support
strong claims about frontier-scale systems. Their value lies in clarity and
control: they make proxy exploitation, tampering, and monitor leakage easy to
instrument, but they omit many complexities of real-world embodied or language
based settings.

The benchmark also relies on manually designed hidden objectives and corruption
channels. That design choice is useful for clean diagnosis, yet it can bias the
space of failures toward what the benchmark author considered important. Future
extensions should therefore be evaluated not only by whether they enlarge the
suite, but by whether they preserve the explicit task schema and artifact
discipline introduced here.

Monitorability results in particular should be interpreted cautiously. A monitor
that performs well in a symbolic benchmark with carefully filtered features does
not imply a practical oversight system for realistic deployments. Conversely, a
poor monitor result may reflect an intentionally restricted audit channel rather
than an impossibility result about external detection in general.

Finally, the local-first execution policy imposes real limits. Training budgets,
seed counts, and model classes are intentionally conservative to keep the repo
reproducible on ordinary hardware. That improves accessibility, but it means the
framework prioritizes breadth of diagnosis over exhaustive scaling studies.
